% Generated draft: requires mathematical and source review.
\documentclass[11pt]{article}
\usepackage[margin=25mm]{geometry}
\usepackage{fontspec}
\setmainfont{Latin Modern Roman}
\usepackage{amsmath,amssymb,mathtools,unicode-math}
\setmathfont{Latin Modern Math}
\usepackage{enumitem,needspace,xurl,hyperref}
\setlength{\parindent}{0pt}
\setlength{\parskip}{5pt}
\setlength{\emergencystretch}{3em}
\newcommand{\N}{\mathbb{N}}
\newcommand{\Z}{\mathbb{Z}}
\newcommand{\Q}{\mathbb{Q}}
\newcommand{\R}{\mathbb{R}}
\newcommand{\C}{\mathbb{C}}
\newcommand{\F}{\mathbb{F}}
\newcommand{\A}{\mathbb{A}}
\newcommand{\PP}{\mathbb{P}}
\newcommand{\E}{\mathbb{E}}
\newcommand{\eps}{\varepsilon}
\newcommand{\dd}{\,\mathrm{d}}
\begin{document}
\section*{UnsolvedMath ID: 10800015}
\textbf{Source problem code:} catalogue entry
\subsection{Definitions and mathematical statement}
Let \(F\) be the real versal deformation of the germ \(f\) intended in the source, with base \(B\), and let \(\mathcal M\subset B\) be its real Morsification parameters. In the usual local interpretation, \(F_u\) at such a parameter has only nondegenerate critical points in the prescribed neighborhood; a real critical point is Morse when its Hessian is nonsingular. Any further distinct-critical-value or complex-critical-point requirements must follow the missing source conventions. Let \(c(u)\) be the source's complexity and \(\rho(f)\) its specified optimal/target complexity, neither of which is defined in the excerpt.

The question is whether
\[\forall u_0\in\mathcal M,\ \exists u_1\in\mathcal M,\ \exists\gamma:[0,1]\to B:\quad
\gamma(0)=u_0,\ \gamma(1)=u_1,\ c(u_1)=\rho(f),\]
with \(\gamma\) generic and every Morse surgery of type \([A_2]\), traversed in the direction \(0\to1\), decreasing the number of real critical points. For an ordinary real birth/death \(A_2\) event, two real Morse critical points are created or annihilated. Thus, writing \(r(u)\) for the number in the fixed local neighborhood, the required direction is
\[r(\gamma(t_*+\eta))=r(\gamma(t_*-\eta))-2\]
for sufficiently small \(\eta>0\) around each such event \(t_*\). Other generic events, if allowed, must obey the source's path conventions; the question is not simply existence of an arbitrary path between chambers.

The formula does not assume that \(c(u)=r(u)\), nor define \(\rho(f)\) as a minimum of \(r\), unless that identification is confirmed from the preceding text. The germ class, base neighborhood, generic-event list, critical-point counting domain and boundary control are essential to the statement and missing here. In particular the directionality forbids inserting births at \(A_2\) events even if a later cancellation lowers the final count.
\subsection{Short English statement}
Can every real Morsification be joined to one of the target complexity by a generic versal-deformation path using only deaths, never births, at its \(A_2\) events? The source omits the definitions of complexity and \(\rho(f)\).
\subsection{Sources}
\begin{itemize}
\item Local UnsolvedMath registry, permanent ID \(10800015\), exact source code \texttt{catalogue entry}.
\item Victor Vassiliev, \emph{A Few Problems on Monodromy and Discriminants}, Problem 6A: \url{https://amj.math.stonybrook.edu/html-articles/Files-2015-2024/15-11/}. Registry access 2026-07-29, html-structure extraction.
\item The summary attributes complete component classifications for simple and parabolic discriminant complements to Vassiliev (2021, 2025), without full citations; it explicitly says those classifications do not settle this path question.
\end{itemize}
\end{document}
